\section{Conclusion and Future Work}
\label{sec:conclusion}

We measured where FANET routing protocols lose packets and tested whether predicting link lifetimes avoids
those losses. On UavNetSim-v1, collisions are almost never the final cause of a loss: the MAC's five attempts
absorb them, and congestion shows up as queueing and expiry instead. Link breaks are the dominant loss for
Greedy forwarding and grow with speed, from 4.1\,\% to 16.0\,\% of generated packets between 5 and 20\,m/s.

LEAP-GR forwards like Greedy unless the chosen link is predicted to expire within $\tauth=0.25$\,s. With the
same MAC and seeds it delivers up to 14.7 points more than Greedy and up to 87 points more than GRAd, for
128 bits per hello and 0.2--0.3 more hops per packet. The ablation shows what the prediction itself is worth:
against a control that switches neighbour after every failed attempt, LEAP-GR gains 0.9 to 3.0 points
at 15--20\,m/s and at $N=40$--50, where breaks are frequent and safe alternatives exist, and loses at 5\,m/s
and under saturation, where its longer paths cost more than the breaks it avoids. The mechanism is mostly
short-horizon: 98.4\,\% of its switches avoid a neighbour that, extrapolated from its last hello, has already
left. Predictive repair after a failure adds nothing once selection is predictive, and weighting neighbours by
their recent retransmission rate lowers retransmissions without improving delivery.

\textbf{Limits.} The study uses a free-space channel, Gauss--Markov mobility, 30\,s runs and three seeds per
configuration; the predictor assumes that each UAV knows its own velocity and that destination positions are
known, as for Greedy.

\textbf{Future work.} Natural next steps are a predictor that tolerates noisy positions and fading (a Kalman
filter, as in~\cite{cui2022tarraq}), coupling the prediction with the platform's virtual-force topology
control to keep links alive rather than only avoiding dying ones, and an evaluation on a ray-traced urban
channel.
